\section{Recomendaciones de despliegue y depuración del sistema}

Para desplegar este Compound AI completo en un producto real, es necesario establecer interfaces claras entre search, APAC y DSPy, y agregar monitoreo a cada resultado intermedio:
\begin{itemize}
    \item El Search trace output debe registrar los cost/constraint flags, para facilitar el debug;
    \item El multi-step question log de APAC debe conservar el intent del usuario más la agent ask list;
    \item El latent cost \(C\) y el solver output de DSPy deben reportarse juntos, para poder alinear distintos benchmark.
\end{itemize}

\begin{importantbox}{Puntos clave de depuración y despliegue}
1) Agregar monitor: en inference, establecer un sanity check para cada trace, question y plan; 2) agregar caché: el search trace suele repetirse, por lo que puede almacenarse en caché y reutilizarse; 3) agregar fallback: cuando el agent judge señale una violación de restricción, retroceder a la capa del hybrid solver para replanificar; 4) agregar control de versiones: cada entrenamiento DPO debe ir acompañado de la versión del modelo judge.
\end{importantbox}

\begin{knowledgebox}{Cómo mantener transparente la cadena}
Escribir un structured log (trace id, agent response, solver plan) en cada etapa del pipeline, y proporcionar para cada log un mapeo de \texttt{trace id -> user request}; de este modo, incluso si el sistema falla en producción, es posible reproducir con rapidez la causa del error.
\end{knowledgebox}

\subsection{Resumen de la sección}

Al desplegar Compound AI, es necesario que la salida de cada módulo sea observable y reproducible, y establecer un enlace de versiones entre judge, solver y trace, para poder ubicar con rapidez las fallas de razonamiento y ejecutar reversiones.

